% ECONOMICS_PROBLEM: {"id": "J62-418", "title": "Causal channels of family advantage", "jel_code": "J62", "source_id": "OP-0363"}
% !TeX program = xelatex
\documentclass[11pt,a4paper]{article}
\usepackage[margin=23mm]{geometry}
\usepackage{fontspec}
\setmainfont{Latin Modern Roman}
\usepackage{amsmath,amssymb,mathtools}
\usepackage{xcolor,enumitem,booktabs,longtable,array,xurl,hyperref}
\definecolor{muted}{HTML}{53636C}
\hypersetup{colorlinks=true,urlcolor=blue,linkcolor=blue}
\setlength{\parindent}{0pt}
\setlength{\parskip}{5pt}
\newcommand{\R}{\mathbb R}
\newcommand{\E}{\mathbb E}
\newcommand{\N}{\mathbb N}
\newcommand{\ind}{\mathbf 1}
\newcommand{\Var}{\operatorname{Var}}
\newcommand{\Cov}{\operatorname{Cov}}
\newcommand{\supp}{\operatorname{supp}}
\DeclareMathOperator*{\argmin}{arg\,min}
\DeclareMathOperator*{\argmax}{arg\,max}
\newcommand{\entrymeta}[3]{{\sffamily\small\color{muted}\textbf{\texttt{#1}}\quad|\quad #2\par #3\par}}
\newcommand{\Problemref}[1]{\texttt{#1}}
\newcommand{\JELref}[2]{\texttt{#2} (JEL \texttt{#1})}
\begin{document}
\section*{Causal channels of family advantage}
\textbf{Problem ID:} \texttt{J62-418}\quad \textbf{JEL:} \texttt{J62}\par
% BEGIN ECONOMICS STATEMENT
\label{problem:OP-0363}
{\sffamily\small\color{muted}Original subject group: Inequality and economic mobility\par}
\label{definitions:problem:30007645}
\textbf{Audit classification:} Background; proposed extension. Author institutional abstract and NBER metadata checked; NBER full text returned 403. The abstract reviews causal approaches but does not explicitly certify the catalogue decomposition as open.
\subsection*{Definitions and precise research target}
\textbf{Complete editorial problem.} Fix a birth cohort, a parental-income range, and a measure \(Y_i\) of adult economic status. Let \(R_i\) denote parental financial resources during childhood and \(H_i\) parental skills, preferences, and information. Let \(N_i\) represent family networks and \(E_i\) the child's institutional environment. Potential outcomes \(Y_i(r,h,n,e)\) are defined only for declared feasible interventions; they do not imply that every inherited characteristic can be manipulated. Holding prereform \(H_i,N_i,E_i\) fixed defines a controlled resource contrast rather than a total effect that includes their responses. The resource effect is \(\tau_R(r,r';x)=E[Y_i(r,H_i,N_i,E_i)-Y_i(r',H_i,N_i,E_i)\mid X_i=x]\), where \(X_i\) is measured before intervention. Analogous contrasts define network or information interventions. Identify these effects using credible income shocks, randomized supports, adoption or family comparisons with explicit assumptions, and longitudinal measurement. Determine how parental resources interact with skills and constraints rather than assigning all parent--child correlation to money. Mediation claims require additional assumptions beyond total-effect identification. Specify exposure timing, selection into parenthood, and whether parental responses are included. The proposed research task is a context-specific causal decomposition and its policy implications; the cited reviews are not evidence that all mechanisms lack established effects or that one universal decomposition exists.

\textbf{Definitions and admissible scope.} Define interventions as actual resource, information, or network programs during specified ages, with consistency across recipients and a stated allocation of spillovers. A controlled-resource outcome \(Y_i(r,h,n,e)\) is defined only where holding the other inputs fixed is technologically feasible. The ordinary total effect of an income transfer allows parental time, information, neighborhood, and networks to respond. Those are different estimands. Baseline covariates \(X_i\) precede exposure, adult status is measured at a fixed age, and nonemployment has an explicit value. Adoption comparisons require placement-selection assumptions and do not randomly assign genetic traits. A decomposition into channels is identified only if the required joint interventions and mediation restrictions are supplied; family correlations alone do not produce causal percentages. Fix admissible intervention domains for \(r,h,n,e\), and require all controlled combinations used in a contrast to be jointly feasible for the same baseline cohort. Let the adult measurement age, welfare/status scale, and baseline covariate law be fixed and outcomes integrable. The required observational or experimental law and its selection, interference, and measurement restrictions are inputs. A controlled network contrast, for example, is \(E[Y_i(r,h,n,e)-Y_i(r,h,n',e)\mid X_i=x]\) on common support; total effects of feasible family programs instead allow the other inputs to respond. These are separate targets. If the admissible law leaves them underidentified, characterize the identified set rather than assigning residual parent--child association to a channel.
\subsection*{Variants and assumption changes}
\begin{enumerate}
\item Resource-by-age intervention (editorial): assign equal present-value transfers during early childhood or adolescence and compare adult outcomes, allowing parental responses. This isolates timing of a feasible policy rather than a genetic channel.
\item Network-by-information intervention (editorial): randomize mentoring access and validated application information factorially. Estimate interaction effects at fixed transfer support and state whether the objective is total policy effects or controlled-input effects.
\end{enumerate}
\subsection*{References and source audit}
\textbf{Support and access.} Author institutional abstract and NBER metadata checked; NBER full text returned 403. The abstract reviews causal approaches but does not explicitly certify the catalogue decomposition as open. \textbf{Locator:} BFI abstract, June 3, 2021.
\begin{enumerate}\raggedright
\item\hypertarget{definitions:source:30007645:1}{} Magne Mogstad; Gaute Torsvik. 2021. \emph{Family Background, Neighborhoods and Intergenerational Mobility}. Abstract, March 2022 revision, discussion of family environment, genetics, and causal family characteristics.
\url{https://www.nber.org/papers/w28874}.
DOI: \url{https://doi.org/10.3386/w28874}.
\end{enumerate}

\subsection*{Common conventions: Source mathematical conventions}

These conventions apply to each entry unless explicitly replaced.

\textbf{Models and identification.} A primitive model \(m\) specifies all probability laws, feasible choices, information, equilibrium and measurement rules used by its target. The admissible class \(\mathcal M\) is fixed independently of outcomes. If \(P_O(m)\) is its observable probability law and \(\tau(m)\) the target, its identified set at observable law \(P\) is
\[
 \mathcal I_\tau(P)=\{\tau(m):m\in\mathcal M,\ P_O(m)=P\}.
\]
Point identification means that this set is a singleton. An empty set signals incompatibility of data and assumptions. Coordinate bounds need not describe the joint set. Bounds are sharp only if every retained target is attainable or arbitrarily approachable by compatible models. Finite bounds require explicit restrictions on unobserved quantities; unrestricted confounding, hidden wealth, extrapolation or arbitrary equilibrium selection can yield uninformative sets.

\textbf{Causal interventions.} A potential outcome \(Y_i(p)\) is the outcome under a complete policy regime \(p\), including timing, financing, information and market spillovers. The target population and units are fixed before treatment. With interference, use the complete assignment vector or a justified exposure mapping; individual no-interference notation alone cannot identify an economy-wide intervention. A difference of mean potential outcomes does not require the cross-world joint distribution. Conditioning on post-treatment migration, survival, enrollment or employment changes the estimand and needs separate justification.

\textbf{Statistical statements.} An estimator, confidence set, forecast or test specifies a sampling law, model class and loss/coverage criterion. Uniform \((1-\alpha)\)-coverage means \(\inf_{m\in\mathcal M}\Pr_m\{\tau(m)\in C_n\}\geq1-\alpha\), or an explicitly asymptotic version. The whole parameter space trivially has coverage, so informativeness requires a stated width, risk or power objective. A forecasting improvement is relative to a named benchmark, information set and evaluation population.

\textbf{Equilibrium and optimization.} An equilibrium comprises feasible plans, optimality, beliefs where needed, budget constraints and all market/resource clearing. The primitives or a stated benchmark define each correspondence \(\mathcal E(p;m)\). Nonemptiness is assumed or proved, never inferred just from compact policy sets. A selected-equilibrium target uses a declared selection rule; a robust target quantifies over all permitted equilibria. Use suprema or infima unless attainment is established. An \(\varepsilon\)-optimal policy is feasible and within \(\varepsilon\) of the finite optimum value. Report an empty feasible set separately.

\textbf{Welfare and accounting.} Normative utility, interpersonal normalization, social weights, numeraire, discounting and the use of public receipts are primitives. Behavioral utility need not coincide with the welfare criterion. Household welfare includes ownership income and rents once; adding profits again to compensating variation generally double-counts them. A monetary equivalent is defined by an explicit transfer and price convention, with monotonicity and range conditions. Average effects, mechanism contributions, transfers and welfare losses are different targets. Infinite sums require convergence and dynamic feasible plans require terminal or no-Ponzi conditions.

\textbf{Source versions.} A working-paper date, revision date and journal publication year are distinguished. A DOI for a journal version is not automatically the DOI of an earlier manuscript. Locators refer to the stated version and printed pages where specified. A metadata-only check does not verify a claim about a full-text conclusion. Every entry's source subsection narrows the provenance claim accordingly.
% END ECONOMICS STATEMENT
\end{document}
